\chapter{Kiến thức nền tảng}
\label{chap:related_work}

Chương này tóm tắt các khái niệm nền tảng cần thiết cho đề tài: xử lý sự cố mạng, phân tích nguyên nhân gốc rễ, mô hình ngôn ngữ lớn, AI Agent và hướng tiếp cận Self-Evolving AI Agents.

\section{Xử lý sự cố mạng và phân tích nguyên nhân gốc rễ}
\label{sec:network_troubleshooting_rca}

Mạng viễn thông hiện đại bao gồm nhiều lớp thiết bị, giao thức và dịch vụ phụ thuộc chặt chẽ với nhau. Một thay đổi nhỏ trong cấu hình định tuyến, trạng thái liên kết, dịch vụ nền hoặc chính sách điều khiển luồng có thể gây lỗi dây chuyền tới nhiều thành phần. Vì vậy, xử lý sự cố mạng không chỉ dừng ở phát hiện bất thường, mà còn cần định vị thiết bị lỗi và xác định nguyên nhân gốc rễ.

Trong đề tài này, quy trình chẩn đoán được chia thành ba tác vụ chính. Thứ nhất là \textbf{detection}, xác định mạng có đang gặp sự cố hay không. Thứ hai là \textbf{localization}, xác định thiết bị hoặc thành phần bị ảnh hưởng. Thứ ba là \textbf{root cause analysis} (RCA), xác định tên nguyên nhân gốc rễ. Về tổng quát, RCA có thể được xem là bài toán suy luận ngược từ trạng thái mạng, log, cấu hình và triệu chứng quan sát được để tìm nguyên nhân phù hợp nhất:
\begin{equation}
\hat{c} = \underset{c \in C}{\mathrm{argmax}}\; p(c \mid U, Y_t, s_t),
\end{equation}
trong đó $U$ là cấu hình mạng, $Y_t$ là dữ liệu quan sát, $s_t$ là triệu chứng tại thời điểm $t$ và $C$ là tập nguyên nhân có thể xảy ra.

\section{Mô hình ngôn ngữ lớn và AI Agent}
\label{sec:llm_agent}

Mô hình ngôn ngữ lớn (LLM) có khả năng hiểu văn bản kỹ thuật, lập luận nhiều bước, sinh giải thích và sử dụng công cụ thông qua API. Trong bài toán xử lý sự cố mạng, LLM có thể đọc mô tả sự cố, phân tích log, gọi công cụ kiểm tra cấu hình và tổng hợp bằng chứng để đưa ra kết luận. Tuy nhiên, LLM thuần túy dễ gặp các hạn chế như thiếu tri thức miền cụ thể, suy luận thiếu kiểm chứng hoặc gọi công cụ sai tham số.

AI Agent mở rộng LLM bằng cách đặt mô hình vào một vòng lặp quan sát--suy luận--hành động. Agent có thể nhận nhiệm vụ, chọn công cụ phù hợp, thực thi truy vấn trên môi trường mạng, quan sát kết quả và tiếp tục điều tra cho đến khi đủ bằng chứng. Các workflow phổ biến gồm ReAct \cite{react2023}, Plan \& Execute \cite{planexecute2023} và Reflexion \cite{reflexion2023}. Điểm mạnh của agent là kết hợp suy luận ngôn ngữ với dữ liệu thực tế từ công cụ, giúp kết quả chẩn đoán có cơ sở hơn so với câu trả lời sinh trực tiếp.

\section{Self-Evolving AI Agents}
\label{sec:self-evolving-agents}

Self-Evolving AI Agents là các agent có khả năng tự cải thiện từ kinh nghiệm trong quá trình hoạt động. Thay vì dùng một prompt, bộ nhớ hoặc mô tả công cụ cố định, agent có thể cập nhật tri thức sau mỗi ca xử lý dựa trên kết quả đánh giá, phản hồi hoặc vết tương tác. Điều này phù hợp với môi trường mạng viễn thông vì topology, giao thức, dịch vụ và dạng lỗi có thể thay đổi theo thời gian.

Các thành phần thường có thể tiến hóa gồm:
\begin{itemize}
    \item \textbf{Bộ nhớ:} Lưu các quy trình chẩn đoán, bài học hoặc mẫu lỗi đã được kiểm chứng để tái sử dụng trong các incident tương tự.
    \item \textbf{Ngữ cảnh và prompt:} Cải thiện cách đưa thông tin vào mô hình, giúp agent tập trung vào bằng chứng quan trọng.
    \item \textbf{Công cụ:} Cải thiện mô tả, ví dụ sử dụng và điều kiện tiền đề của công cụ để giảm lỗi gọi công cụ.
    \item \textbf{Chiến lược suy luận:} Điều chỉnh thứ tự kiểm tra giả thuyết, cách loại trừ nguyên nhân và cách dừng điều tra.
\end{itemize}

Trong nghiên cứu này, hệ thống tập trung vào hai hướng tiến hóa có rủi ro triển khai thấp hơn so với fine-tuning mô hình: tiến hóa bộ nhớ quy trình và tiến hóa tài liệu công cụ.

\section{Các nghiên cứu liên quan trực tiếp}
\label{sec:related_frameworks}

Đề tài chỉ tập trung vào hai hướng nghiên cứu liên quan trực tiếp đến phần hiện thực. Thứ nhất, Skill-Pro nghiên cứu procedural memory cho LLM agent, trong đó kinh nghiệm xử lý được biểu diễn thành các kỹ năng có thể truy xuất và tái sử dụng trong quá trình suy luận dài \cite{skillpro2024}. Ý tưởng này phù hợp với bài toán xử lý sự cố mạng vì mỗi incident thành công có thể trở thành một quy trình chẩn đoán hữu ích cho các tình huống tương tự.

Thứ hai, DRAFT tập trung vào việc cải thiện khả năng sử dụng công cụ của LLM thông qua tinh chỉnh tài liệu mô tả công cụ \cite{draft2024}. Thay vì thay đổi mã nguồn công cụ, hệ thống học từ các lần gọi công cụ thực tế để bổ sung mô tả, điều kiện sử dụng, ràng buộc tham số và hướng dẫn diễn giải kết quả trả về.

Từ hai hướng này, đề tài xây dựng hai mô-đun tương ứng: \textbf{Procedural Memory} cho tiến hóa bộ nhớ quy trình và \textbf{Tool Refinement} cho tinh chỉnh tài liệu công cụ. Cách tiếp cận này giúp agent cải thiện hành vi mà không cần cập nhật tham số mô hình nền, nhờ đó dễ kiểm soát và phù hợp hơn với yêu cầu an toàn của hệ thống vận hành mạng.
